% SOURCE_AUTHORITY: sga5_fr_workpass.tex, lines 11231--11266 (LF-normalized unit).
% SOURCE_SHA256: 791F4EFFC5E02832D5D77ED03518C8156D6F07E4C8238B03545DB93D883FBB28
\subsection*{9.5}

Demonstremos finalmente como o que precede implica o teorema. Denotemos por $p:X\times\mathbb P^1_k\to X$ a primeira projeção, de modo que $i=p\circ i_1$. Se $y\in C^*(Y)$, tem-se evidentemente
\[
i^*i_*(y)=(i_1)^*(\bar x),\qquad \bar x=p^*i_*(y).
\tag{9.5.1}
\]
Por (9.1.1), deduz-se que
\[
i^*i_*(y)=(g_1)_*\bigl[g_1^*i_1^*(\bar x)c_d(\check E)\bigr],
\]
Isto é,
\[
i^*i_*(y)=(g_1)_*\bigl[j_1^*f_1^*(\bar x)\cdot c_d(\check E)\bigr].
\tag{9.5.2}
\]
Para demonstrar o teorema, pode-se, por linearidade, supor que $y$ é o ciclo irredutível associado a uma subvariedade $T$ de $Y$, de modo que $\bar x$ corresponde ao ciclo $T\times\mathbb P^1$ de $Z=X\times\mathbb P^1$. Denotando por $\theta$ a transformada estrita por $f_1$ de $T\times\mathbb P^1$, isto é, a aderência da imagem inversa de $T\times\mathbb P^1$ sobre $Z\doteq Y$, tem-se uma relação da forma
\[
f_1^*(\bar x)=\theta+(j_1)_*(w),\qquad w\in C^{\bullet}(\widehat N).
\tag{9.5.3}
\]
Substituamos esta expressão em (9.5.2). Resulta então de 9.4 que
\[
i^*i_*(y)=(g_1)_*\bigl[j_1^*(\theta)c_d(\check E)\bigr],
\tag{9.5.4}
\]
de onde, por (9.1.4),
\[
i^*i_*(y)=s^*j_1^*(\theta).
\tag{9.5.5}
\]
Mas, como o esquema (reduzido) subjacente a $\theta$ é o esquema obtido pela explosão de $T\times\mathbb P^1$ ao longo de $T=(T\times\mathbb P^1)\cap Y$, ele corta transversalmente a seção zero de $\widehat N$ segundo a imagem de $T$ por essa seção zero. Por conseguinte,
\[
j_1^*(\theta)=s_*(y),
\]
que, tendo em conta (9.5.5) e (9.1.3), termina a demonstração do teorema.
